\chapter{总结}\label{chap:conclusion}

这次页表改造的核心，是把 Unix V6++ 的用户内存管理从“连续进程镜像”的思路，推进到“按页框分配、按页表描述地址空间”的思路。实现上主要完成了下面几层变化：

\begin{itemize}
  \item 物理页框由位图管理，\texttt{PageManager} 维护页引用计数；
  \item 每个进程有自己的页目录和私有用户页表；
  \item 页目录中同时保留公共用户页表、内核页表和进程私有页表；
  \item 上下文切换时动态写入 \texttt{CR3}，真正切换地址空间；
  \item \texttt{fork()}、\texttt{exec()}、\texttt{brk()}、栈增长和 \texttt{exit()} 都按页维护映射关系；
  \item \texttt{pageDemo} 验证了同一虚拟地址布局下的父子进程写入隔离。
\end{itemize}

从运行结果来看，现在系统已经具备了一个分页式地址空间应有的基本性质：进程看到的是自己的虚拟地址空间，物理页框由页表决定，父子进程可以拥有相同的虚拟地址，却不会因为一次写入互相破坏数据。

后续如果继续完善，可以把写时复制异常处理链路补完整，让只读共享页在第一次写入时自动分配新页；也可以继续清理旧的连续内存字段，让进程镜像描述彻底转向页表和页框数组。就当前课程设计目标来说，页框分配、私有页表、动态页目录切换和运行验证这几部分已经把核心思想串起来了。
